\documentclass[12pt,letterpaper]{article}
\usepackage[T1]{fontenc}
\usepackage{amsmath,amssymb,newtxtext,newtxmath,microtype,setspace}
\usepackage[margin=1in]{geometry}
\usepackage{graphicx,booktabs,array,float,caption,hyperref}
\hypersetup{hidelinks}
\captionsetup{font=small,labelfont=bf,labelsep=period}
\setstretch{1.5}
\setlength{\parindent}{0.25in}
\setlength{\parskip}{0pt}
\setlength{\emergencystretch}{1em}
\renewcommand{\thesection}{\Roman{section}}
\renewcommand{\thesubsection}{\Alph{subsection}}
\renewcommand{\thetable}{\Roman{table}}
\setcounter{section}{5}
\setcounter{figure}{2}
\setcounter{page}{28}
\newcommand{\C}{\widehat C}
\newcommand{\SR}{\operatorname{SR}}
\newcommand{\figpage}[3]{\clearpage\begin{figure}[p]\centering
\includegraphics[width=\textwidth,height=.84\textheight,keepaspectratio]{#1.pdf}
\caption{#2}\label{#3}\end{figure}\clearpage}
\graphicspath{{../figures/}{../../empirical_three_experiments_20261009/figures/}}
\begin{document}
\section{Empirical Portfolio Learnability}

Section V distinguishes the opportunities represented by a portfolio policy class from those that can be learned from finite return histories. We investigate this distinction in three empirical experiments. The first holds the Gaussian representation fixed and traces performance as regularization activates additional managed-portfolio directions. The second changes the amount of available history on common evaluation dates. The third reconstructs actual stock holdings to identify the investment strategies, risk exposures, and incremental portfolio value represented by different parts of the managed-payoff spectrum.

We retain the existing U.S. monthly stock panel of Jensen et al.\ (2023) and the characteristic preprocessing following Didisheim et al.\ (2024). Formation dates cover January 1963 through December 2024. Nano stocks are excluded; 130 characteristics are chosen from the 153-characteristic dictionary using coverage in the initial 1963--1972 period. Stock-months with more than 30\% missing retained characteristics are removed. Monthly cross-sectional ranks lie in $[-0.5,0.5]$, and remaining missing characteristic values receive zero. The managed payoff is the cross-sectional average of characteristic-based feature payoffs, so a fitted coefficient vector defines the policy $W$ through the existing $1/N_t$ stock-weight convention.

We accept the retrospectively constructed JKP characteristic library as the academic information set for this analysis. All results remain conditional on the existing retained, complete-payoff panel. The economic interpretation preserves that universe and its preprocessing; it does not introduce a new information set.

The Gaussian map uses 10,000 fixed Random Fourier Features, seed zero, and the existing initial-sample bandwidth of approximately 4.3921. The 120 positive candidate penalties are frozen using the initial managed-payoff spectrum. We preserve annual refitting and monthly characteristic-dependent weights. In the expanding-history experiments, the previous five formation years provide chronological validation under response-one loss; the policy is then refitted using all available training and validation payoffs. A January formation decision uses returns realized through that January and is evaluated from February through the following January. Thus the 47 decision years 1978--2024 yield 564 returns from February 1978 through January 2025. All period labels below denote decision years.

Sharpe ratios are annualized as $\sqrt{12}$ times the monthly mean divided by its sample standard deviation. Response-one loss is evaluated on the original, unscaled managed payoff. Economic mean returns and volatility use the existing common coefficient scale $\kappa=0.0374000$, frozen at the first decision date; this scaling does not affect Sharpe ratios. Results are before trading and borrowing costs. We do not impose a new leverage cap or ex post volatility normalization.

Uncertainty bands use 5,000 paired circular block-bootstrap resamples of the realized monthly policy returns, with 12-month blocks. Annual complexity comparisons use three-year blocks. These are pointwise, conditional intervals: the entire training and selection process is not re-estimated within each resample. Blocks of 6 and 24 months, and 2 and 5 years for the main complexity contrast, provide sensitivity checks. The five-year final period contains limited independent history, and overlapping estimation windows and changing economic conditions further restrict inference. Neither the period comparisons nor the spectral comparisons are adjusted for multiple testing.

\subsection{Profitable Active Complexity Across Test Periods}

For each fixed penalty $\lambda$, we reconstruct the complete return path by refitting every year. The horizontal coordinate in Figure~\ref{fig:complexity} is the arithmetic mean of $\C_y(\lambda)$ across the annual refits in that test block; every refit contributes twelve evaluation months. The in-sample curve is the mean of the corresponding annual refit-sample Sharpe ratios. These training samples overlap, so this curve is a fit diagnostic rather than the performance of a distinct historical trading strategy. The axes are common across periods; in-sample Sharpe is displayed logarithmically because very weak regularization approaches interpolation.

The policy selected by prior validation can change $\lambda$ each year. Its mean complexity is marked by a vertical line, but its realized Sharpe is reported separately in Table~\ref{tab:selected}; it is not a point on a fixed-penalty performance curve. The crosses in Figure~\ref{fig:complexity} identify descriptive maxima on the test-period grid. Their locations are selected with hindsight and have no investable interpretation.

\begin{table}[htbp]
\centering\small\setstretch{1.12}
\caption{Chronologically selected Gaussian policy, reported separately from fixed-penalty paths}\label{tab:selected}
\begin{tabular}{lrrrrr}
\toprule
Decision years & Months & Mean $\C$ & Sharpe & 95\% interval & Loss \\
\midrule
1978--1989 & 144 & 55.6 & 5.50 & [4.80, 6.86] & 0.288 \\
1990--1999 & 120 & 156.2 & 5.81 & [4.85, 7.35] & 0.283 \\
2000--2009 & 120 & 152.1 & 2.81 & [2.01, 3.95] & 0.756 \\
2010--2019 & 120 & 187.3 & 2.98 & [2.46, 3.69] & 0.604 \\
2020--2024 & 60 & 45.9 & 0.76 & [0.21, 2.43] & 1.276 \\
\bottomrule
\end{tabular}
\par\vspace{5pt}
\begin{minipage}{\textwidth}\footnotesize Each decision year runs from February through the following January. The selected penalty can change across annual decisions. Complexity is averaged over those refits, and Sharpe is calculated on the concatenated OOS returns. All intervals are pointwise.\end{minipage}
\end{table}


The descriptive maximizing complexities are approximately 154, 145, 131, 288, and 16 across the five periods. The early periods display large gains from moving away from nearly inactive policies, followed by comparatively modest deterioration at the most weakly regularized end of the grid. In 2000--2009 the curve is flatter, while the 2010--2019 maximum occurs farther into the spectrum. In 2020--2024, the maximum occurs much earlier and the uncertainty band is wide. These are different economic environments rather than estimates of one universal optimal complexity.

Historical fit continues to improve far beyond the range producing material out-of-sample gains. Yet the data do not show that every additional direction is harmful: even the highest-complexity grid endpoint outperforms the lowest-complexity endpoint in the first four blocks. The chronologically selected policy's Sharpe falls from 5.50 and 5.81 in the first two blocks to 2.81, 2.98, and 0.76 subsequently. The fixed representation therefore supports substantially different amounts of useful active complexity across periods. These findings support regularization as a restriction on what can be learned, while leaving substantial uncertainty about the precise location of an economically relevant peak.

\figpage{E1_complexity_by_period}{Fixed-penalty Gaussian performance paths in five disjoint test periods. Left: out-of-sample Sharpe and pointwise 95\% block-bootstrap bands. Right: mean annual refit-sample Sharpe. Dashed lines mark mean validation-selected complexity only; Table~\ref{tab:selected} reports that policy's separate performance.}{fig:complexity}

\subsection{Market History and Learnable Complexity}

We next compare trailing histories $T\in\{60,120,240,360\}$ months. The Gaussian features, bandwidth, candidate penalties, scaling, and evaluation dates are identical. To isolate history length, every specification reserves its last 20 months for chronological validation, estimates candidates on the preceding $T-20$ months, and refits the selected policy on all $T$ observations. This common validation window is an explicit change from the five-year expanding-history protocol: a five-year holdout would leave no training sample when $T=60$. It is held fixed across the four histories rather than being allowed to vary with $T$.

The common sample comprises 32 January decisions, 1993--2024, and 384 payoffs from February 1993 through January 2025. Table~\ref{tab:history} reports mean effective complexity, median selected penalty, Sharpe, and response-one loss. Mean complexity rises from 34.4 to 59.9, 125.8, and 170.1 as history increases. The paired $T=360$ minus $T=60$ complexity difference is 135.7, with a three-year block interval of $[94.2,177.0]$, and is positive in 28 of 32 refits. A within-year regression of complexity on $\log T$ gives a mean slope of 76.7, with interval $[53.3,100.2]$. These statistics measure activation under the fixed empirical selection procedure; they are not estimates of a population complexity law or its exponent.

\begin{table}[htbp]
\centering\small\setstretch{1.12}
\caption{Common-date history comparison, February 1993--January 2025}\label{tab:history}
\begin{tabular}{rrrrr}
\toprule
$T$ & Mean $\C_T$ & Median $\widehat\lambda_T$ & OOS Sharpe & Response-one loss \\
\midrule
60 & 34.4 & $30.72\times10^{-9}$ & 2.367 & 0.778 \\
120 & 59.9 & $30.49\times10^{-9}$ & 2.488 & 0.757 \\
240 & 125.8 & $5.37\times10^{-9}$ & 2.750 & 0.696 \\
360 & 170.1 & $5.18\times10^{-9}$ & 2.629 & 0.727 \\
\bottomrule
\end{tabular}
\par\vspace{5pt}
\begin{minipage}{\textwidth}\footnotesize All specifications have 384 identical evaluation months and a common 20-month validation window. Paired full-period and subperiod contrasts, interval endpoints, annual penalties, and grid-boundary frequencies accompany the replication files.\end{minipage}
\end{table}


More active directions do not translate mechanically into better portfolio performance. Full-period Sharpe increases through $T=240$ and then falls from 2.75 to 2.63 at $T=360$. The latter change is $-0.12$, with paired interval $[-0.44,0.19]$; response-one loss rises by 0.031, with interval $[-0.058,0.119]$. Thus the full-period evidence does not establish a performance advantage for the longest history. The candidate grid also binds in 28\%, 19\%, 16\%, and 16\% of refits, respectively. Selected complexity is therefore conditional on this frozen grid, not an unconstrained optimum.

Figure~\ref{fig:history} separates historical subperiods. In 1993--1999, extending history from 240 to 360 months lowers Sharpe by 1.49 and raises loss by 0.085; the respective paired intervals are $[-2.73,-0.49]$ and $[0.018,0.170]$. In 2010--2019, moving from 60 to 240 months instead raises Sharpe by 1.09 and lowers loss by 0.262, with intervals $[0.37,1.81]$ and $[-0.472,-0.096]$. Both patterns are conditional, exploratory comparisons. In 2020--2024, all four response-one losses exceed one, the loss of a zero-payoff policy, even though additional history increases the selected number of active directions. Rank activation and reliable economic value are distinct.

Local concentration is measured by the top-ten eigenvalue share and the participation rank, $(\sum_j\widehat\mu_j)^2/\sum_j\widehat\mu_j^2$, of each trailing-history uncentered managed-payoff operator. Figure~\ref{fig:local} shows time variation in the top-ten share, although it remains above 99\% throughout. The average participation rank is close to 1.04, illustrating how a dominant second-moment direction can coexist with many weak directions activated by small penalties. Across refits, the correlation between the $T=60$ top-ten share and $\C_{360}-\C_{60}$ is only $-0.049$. This diagnostic does not explain when a longer history is most useful.

The evidence is consistent with a statistical benefit from additional observations and a cost from incorporating less relevant historical observations, but these effects cannot be separately identified from the window comparison. Changes in selected penalties, estimated means, and the economic spectrum occur together. We therefore do not attribute the deterioration in a particular block exclusively to structural change. Nor do we estimate a global spectral exponent or apply a polynomial-decay formula to the Gaussian kernel. More history generally activates more directions in this sample; neither a monotone performance gain nor a stable economic condition guaranteeing that gain is established.

\figpage{E2_history_comparison}{History length and selected active complexity on common evaluation dates. Each column is a disjoint decision-year block. The top row reports average selected complexity; the bottom row reports realized Sharpe. Table~\ref{tab:history} supplies penalties and response-one loss for the common full sample, and the accompanying data tables report all subperiod values.}{fig:history}
\figpage{E2_local_spectrum}{Local concentration of the managed-payoff operator and the complexity gain from more history. Concentration is computed before each decision. The scatter compares two outcomes of the same annual window experiment and is descriptive, not a causal or predictive regime classification.}{fig:local}

\subsection{Investment Strategies in the Managed-Payoff Spectrum}

We now ask what is bought and sold when the existing ridge policy activates weaker managed-portfolio directions. We retain the four training spectral rank groups, 0--10\%, 10--50\%, 50--90\%, and 90--100\%, and the annually validation-selected penalty. Numerically positive eigenvalues exceed $10^{-12}$ times the largest eigenvalue; near ties remain together at the existing boundaries. There is no new portfolio-learning procedure or selection of characteristics. All results in this subsection concern actual ridge policy contributions, at the same common scale, rather than equal-risk diagnostic eigenportfolios.

For each formation month, we evaluate each group's contribution to $W$ on every stock's contemporaneous characteristics and retain its signed weight under the original $1/N_t$ convention. The four stock weights sum to the original fitted stock weight, and their next-month payoffs sum to the original policy payoff. This reconstruction covers 2,198,755 stock-months and 47 annual refits. The accompanying private files identify every position by its verified CRSP PERMNO, date, group, and signed weight. Long and short lists are ordered by position size, never subsequent performance. Company-name mappings are unavailable in the retained extract; identifiers are not replaced by guessed names.

For interpretation, we compute the holding-weighted mean of each input characteristic separately on the long and absolute-short legs. Their difference measures a characteristic tilt independently of gross exposure; we also retain the actual signed notional exposure. Multiplying the difference by 100 expresses it in input percentile points. Family scores average economically related ranked characteristics, with signs chosen to mean more of the named attribute, not higher expected returns. The dictionary covers value, profitability, investment, momentum, recent return, size, liquidity, volatility, equity issuance, accruals, quality, leverage, and seasonality. A financing follow-up, motivated by the largest observed holdings exposures, averages financial-liability growth and net debt issuance. It is explicitly descriptive, not a preregistered hypothesis or an additional trading rule. Definitions and all 130 characteristic exposures accompany the results.\footnote{Economic definitions follow the \href{https://jkpfactors-data.s3.amazonaws.com/documents/Documentation.pdf}{JKP documentation} and its official characteristic-name dictionary. Only definitions and names are used for this interpretation.}

\paragraph{What the four groups hold.}
Figure~\ref{fig:holdings} and Table~\ref{tab:dominant} describe economically distinct, but overlapping, sets of positions. The strongest group favors cash profitability and quality, momentum, and value, with lower accruals, investment growth, volatility, and recent returns. Its long-minus-short family differences are 23.9, 23.5, 22.8, and 17.4 percentile points for profitability, quality, momentum, and value, respectively. Negative recent-return exposure is consistent with short-horizon reversal, alongside positive medium-horizon momentum. Cash-based operating profitability and operating cash flow are its strongest individual positive tilts. This is a broad combination of familiar strategies, not one uniquely identified factor.

Ranks 10--50\% instead counterweight part of that strong-group profile. The corresponding differences are $-7.6$ points for profitability, $-8.6$ for quality, $-5.2$ for momentum, and $+5.4$ for investment. Financial-liability growth is the strongest positive individual exposure, at $+16.2$ points; change in net financial assets is the strongest negative exposure, at $-17.4$. The debt-financing family difference is $+12.1$ points. Equity issuance tilts negative, so liability-financed expansion and share issuance are not interchangeable descriptions. The appropriate interpretation is a financing and investment tilt combined with corrections to the leading group's quality, momentum, and value positions. Calling it another conventional value or quality strategy would misdescribe its holdings.

The 50--90\% group favors higher accruals and lower cash profitability and quality, with slightly negative momentum. Its largest positive tilt is total accruals relative to income; its negative tilts include operating cash flow, the financial-strength score, and quality profitability. The final group has smaller, mixed exposures: relatively more asset liquidity and inventory growth, and less earnings variability. Its broad volatility tilt is negative, but that does not establish a pure low-volatility strategy. Neither group warrants an invariant economic label attached to its individual eigenvectors.

\begin{table}[p]\centering\small\setstretch{1.12}
\caption{Dominant long and short characteristic exposures in actual holdings}\label{tab:dominant}
\begin{tabular}{lp{3.05in}rrr}
\toprule
Ranks & Characteristic & Long & Short & L--S \\
\midrule
0--10\% & Cash-based operating profits-to-book assets & 68.3 & 30.7 & +37.7 \\
0--10\% & Operating cash flow to assets & 68.5 & 31.2 & +37.2 \\
0--10\% & Operating accruals & 35.9 & 63.4 & -27.5 \\
0--10\% & Percent operating accruals & 37.1 & 63.7 & -26.6 \\
10--50\% & Change in financial liabilities & 58.0 & 41.9 & +16.2 \\
10--50\% & Change in noncurrent operating liabilities & 55.3 & 44.3 & +11.0 \\
10--50\% & Change in net financial assets & 41.3 & 58.7 & -17.4 \\
10--50\% & Percent total accruals & 43.1 & 56.3 & -13.2 \\
50--90\% & Percent total accruals & 57.4 & 43.8 & +13.6 \\
50--90\% & Total accruals & 54.8 & 44.9 & +9.9 \\
50--90\% & Operating cash flow to assets & 45.5 & 53.6 & -8.1 \\
50--90\% & Piotroski F-score & 45.8 & 53.2 & -7.4 \\
90--100\% & Liquidity of market assets & 52.7 & 47.0 & +5.8 \\
90--100\% & Inventory change & 52.6 & 46.9 & +5.7 \\
90--100\% & Earnings variability & 45.9 & 55.3 & -9.4 \\
90--100\% & Earnings volatility & 49.0 & 54.2 & -5.1 \\
\bottomrule\end{tabular}\par\vspace{5pt}
\begin{minipage}{\textwidth}\footnotesize Two largest positive and two largest negative differences per group, ranked by average holding exposure, not realized return. Long and short columns show holding-weighted input percentile scores (neutral missing inputs equal 50); L--S is their difference in percentile points. The companion named table reports six exposures on each side for every historical period. Stocks can share multiple characteristics; this is not a disjoint strategy decomposition.\end{minipage}\end{table}

\figpage{EC1_holdings_and_characteristics}{Economic content of actual ridge contributions. Panel A compares long and short characteristic means. Panels B and C show signed notionals and subsequent excess-payoff contributions at the frozen policy scale. Component gross positions are not additive account leverage: holdings across groups can offset on the same stock.}{fig:holdings}

\paragraph{Why little second-moment mass can have financial value.}
Table~\ref{tab:financial} links the spectral quantities to actual positions. Ranks 10--50\% account for only 0.0685\% of training second-moment mass, but their average long and short notionals are 61.48\% and 64.81\% of NAV. Their annual long-leg contribution is 6.00 percentage points and their short-leg contribution is $-3.24$, producing 2.76 points overall, with a pointwise interval of $[1.50,4.07]$. Thus the shorts are not uniformly profitable. Eigenvalue mass measures variation of managed payoffs before application of the selected policy coefficients; it neither measures invested notional nor apportions expected returns. Ridge contributions also depend on estimated means and the selected regularization, so low mass alone does not imply economically negligible positions.

Much of the second group's activity rebalances positions already held by the first. Its 126.29\% component gross raises netted account gross only from 226.68\% to 233.28\%, a 6.60-point increase. Its own annual variance is 0.001547, but twice its covariance with the preceding policy is $-0.001269$; about 82\% of that own variance is offset, leaving a variance increment of 0.000278. The group therefore contributes both expected payoff and diversification, even though its addition slightly increases total volatility.

\begin{table}[p]\centering\small\setstretch{1.12}
\caption{Spectral mass, actual investment scale, and subsequent payoff}\label{tab:financial}
\begin{tabular}{lrrrrr}
\toprule
Ranks & Mass & Gross & Netted gross & Mean & Volatility \\
\midrule
0--10\% & 99.9215 & 226.68 & 226.68 & 24.323 & 8.076 \\
10--50\% & 0.0685 & 126.29 & 233.28 & 2.759 & 3.933 \\
50--90\% & 0.0091 & 32.66 & 232.50 & 0.047 & 0.835 \\
90--100\% & 0.0010 & 7.77 & 232.61 & 0.026 & 0.267 \\
\bottomrule\end{tabular}\par\vspace{5pt}
\begin{minipage}{\textwidth}\footnotesize All entries are percentages. Mass is the average training second-moment share. Gross refers to the individual group; netted gross sums stock weights across all groups through the stated rank boundary before taking absolute values. Mean and volatility are annualized excess-payoff contributions of the individual group, not the cumulative portfolio. No rescaling is introduced.\end{minipage}\end{table}


Contemporaneous characteristic terciles partition each month's actual holdings. Every cell is retained; each family separately accounts for the entire group. Low-profitability stocks contribute 2.04 annual percentage points, low-momentum stocks 1.86, and recent losers 1.99. These sets overlap and their payoffs must not be added. In the financing view, low, middle, and high financing terciles contribute $-0.11$, 1.42, and 1.45 points.

Figure~\ref{fig:joint} examines joint holdings. Relative signed weight density divides a cell's mean stock weight by the group's mean absolute stock weight. We project these nine formation-date cell densities onto additive row and column effects, weighting by contemporaneous stock counts. The residual describes nonadditivity in holdings, not a learned return model. For example, high financing and low profitability have positive density 0.25, whereas low financing and high profitability have density $-0.50$. Their annual payoff contributions are 0.62 and $-0.60$ percentage points, respectively. Nonadditive residuals of approximately $-0.052$ for high financing/low profitability and $+0.067$ for high financing/high profitability temper the simple additive anti-profitability description. These conditional patterns do not establish an interaction premium.

\figpage{EC3_group2_joint_strategies}{Joint characteristic content of ranks 10--50\%. Columns report signed relative weight density, its residual from a formation-date additive row/column fit, and subsequent annualized payoff contribution. Every cell is retained, and colors share a scale within each column. The financing views follow from holdings exposures; no profitable cell is used to define a new strategy.}{fig:joint}

\paragraph{Established risks and incremental out-of-sample value.}
We regress actual contributions on market, size, value, profitability, investment, and momentum factors. The strongest group loads positively on value (0.161), profitability (0.313), and momentum (0.088), with twelve-lag HAC intervals excluding zero. The second group has a negative value loading of $-0.104$, with interval $[-0.194,-0.013]$; its other five loading intervals include zero. The factors explain 17.9\%, 7.1\%, 1.8\%, and 4.3\% of the four groups' return variation. The tables provide all loadings and cell-level factor attributions.

For the second group, the annualized descriptive regression intercept is 3.55\%, with HAC interval $[1.94,5.16]$, exceeding its 2.76\% mean payoff: conventional factor exposures account for approximately $-0.79$ percentage points of the mean. The existing hedge with slopes estimated before each test year leaves 3.60\% mean annual payoff, with interval $[2.28,4.96]$. The regression and chronological hedge are different objects, but both indicate that these six factor exposures do not account for the positive component mean. They do not identify a risk-free opportunity or rule out other economic risks and implementation costs.

Crucially, a positive component mean does not establish a precisely measured improvement to the full investment policy. Adding ranks 10--50\% raises nested Sharpe from 3.012 to 3.284 and lowers response-one loss by 0.0360, but the paired loss interval $[-0.0941,0.0294]$ includes zero. Ranks 50--90\% add only 0.047 annual percentage points and essentially no loss improvement. The final group contributes 0.026 points, with a standalone mean interval including zero, while lowering variance and loss modestly; its incremental loss is $-0.0031$, with an unadjusted pointwise interval $[-0.0071,-0.0001]$.

Figure~\ref{fig:accounting} allocates incremental loss and variance to economic cells, retaining all cross-covariances. For cell payoff $q_b$ within added group $q$, and preceding payoff $p$, its response-one loss allocation is the average of $-2(1-p)q_b+q_bq$, using original unscaled payoffs. These allocations sum exactly to the group's loss change. The corresponding annual variance allocation is $12[2\operatorname{Cov}(p,q_b)+\operatorname{Cov}(q,q_b)]$, using scaled excess payoffs. This is an accounting allocation, not a causal marginal value for independently removing a cell. High-financing stocks allocate $-0.0518$ to loss and $-0.000406$ to annual variance, while low-financing stocks allocate $+0.0361$ and $+0.000571$. Return contribution and diversification contribution can therefore point in different directions. Adding values across alternative characteristic partitions would double-count the same holdings.

\figpage{EC4_group2_incremental_attribution}{Economic allocation of the second group's incremental value. Each family separately sums to the same group mean, loss change, and variance change. Intervals are pointwise paired block intervals conditional on fitted policy paths. Negative loss allocations improve response-one fit; negative variance allocations reduce risk after accounting for covariance with preceding groups.}{fig:accounting}

\paragraph{Stability of the investment interpretation.}
Figure~\ref{fig:styles} repeats the holdings interpretation in the five existing historical periods. The first-to-last cosine of the 13-family exposure vector is 0.914 for ranks 0--10\%, with all family signs preserved. It is 0.700 for ranks 10--50\%, $-0.186$ for ranks 50--90\%, and 0.443 for ranks 90--100\%. These are economic-exposure similarities, distinct from correlations of realized returns or adjacent coefficient vectors. Rank groups are not fixed sets of firms or immutable individual directions.

The second group's investment tilt remains positive and its profitability and momentum tilts negative in all five period averages. Its value tilt changes sign, however, and the financing difference declines from 25.0 points in 1978--1989 to 2.3 in 2000--2009 before recovering to 7.5 in 2020--2024. Its average component gross falls from 178.7\% in the 1990s to 33.4\% in the final period. Thus both economic composition and investment intensity change. The financing interpretation is strongest early in the sample; it is not a constant loading assumed to hold across decades. The much weaker stability of ranks 50--90\% cautions especially against naming its eigenvectors as permanent strategies.

The investment opportunities learned as regularization admits weaker directions are consequently best described in layers. Strong directions implement a stable combination of cash profitability, quality, momentum, and value. The economically meaningful next group makes financing-related and conditional stock-weight adjustments that partly offset those dominant positions, adding expected payoff with substantial covariance offsets. The deepest groups provide much smaller and less stable mean contributions, with some modest diversification value. This evidence supports learning refinements to an existing multivariate investment policy, rather than interpreting every weak eigenvector as a new, separately identified anomaly.

\figpage{EC2_economic_stability}{Historical stability of group-level economic exposures. Each panel uses the same color scale and the existing decision-year periods. Cosines compare the first and final vectors of 13 core family exposures; the separately motivated financing follow-up is displayed but excluded from that comparison. Monthly trajectories, annual characteristic exposures, and period-specific return and factor attributions accompany the figure.}{fig:styles}

\end{document}
